\centerline{\bf \Huge Вокруг ортоцентра}

\begin{flushright}
        \scriptsize{Эпизод 2. Незнакомые потолки}
\end{flushright}

\problem{1}
Предположим, что две вершины треугольника и его описанная окружность зафиксированы, а третья вершина движется по этой описанной окружности. По какой траектории будет двигаться ортоцентр треугольника? (Не забудьте, что вершина может двигаться по двум разным дугам!)

\problem{2}
Докажите, что радиусы описанной окружности, проведенные к вершинам треугольника, перпендикулярны соответствующим сторонам ортотреугольника.

\problem{3}
В остроугольном треугольнике $A B C$ угол при вершине $A$ равен $60^{\circ}$. Известно, что $B C=1$. Найдите длину отрезка $B_1 C_1$, если $B_1$ и $C_1$ это основания высот треугольника, опущенных из вершин $B$ и $C$.

\problem{4}
Докажите, что сумма квадратов расстояния от вершины треугольника до его ортоцентра и длины стороны, противолежащей этой вершине, равна квадрату диаметра описанной окружности.

\problem{5}
Треугольник $A B C$ с ортоцентром $H$ вписан в окружность $\omega$. Окружность, построенная на отрезке $A H$ как на диаметре, пересекает $\omega$ в точке $S \neq A$. Докажите, что прямая $S H$ делит отрезок $B C$ пополам.

\problem{6}
Докажите, что середины трех отрезков, соединяющих ортоцентр с вершинами треугольника, лежат на окружности Эйлера этого треугольника.

\problem{7}
Четырехугольник $A B C D$ вписанный. Докажите, что ортоцентры треугольников $A B C, B C D$, $C D A$ и $D A B$ образуют четырехугольник, равный исходному.

\problem{8}
\textbf{Точка Анти-Штейнера.}
Прямая $\ell$ проходит через ортоцентр треугольника $A B C$. Докажите, что три прямые, симметричные прямой $\ell$ относительно сторон треугольника, пересекаются в одной точке.